%% kapa: exam=1 round=30 year=2019 period=1 form=A subject=real_estate subject_ko=부동산학원론 seq=29 number=109
%% source: Q-Net 감정평가사 기출문제 (cst00309) artlSeq=5205205 file=extracted/1008/5205205/2019년 제30회 감정평가사 1차시험 문제지/2019년 제30회 감정평가사 1차시험 문제지/2019년도 제30회 감정평가사 1차 1교시 A형.hwp
%% conversion: kordoc 4.24.51 -> tools/gen1.py (auto); review=proofread-sonnet; flags=none
\begin{question}{109}
우리나라의 부동산정보 관리정책에 관한 설명으로 옳은 것은?
\begin{choices}
  \item 부동산거래 계약과 신고 등에 관한 정보체계 구축의 법적 근거는 ｢공간정보의 구축 및 관리 등에 관한 법률｣이다.
  \item 국토교통부장관 또는 시장ㆍ군수ㆍ구청장은 정보의 관리를 위하여 관계 행정기관이나 그 밖에 필요한 기관에 필요한 자료를 요청할 수 있으며, 이 경우 관계 행정기관 등은 특별한 사유가 없으면 요청에 따라야 한다.
  \item 광역시장ㆍ도지사는 적절한 부동산정책의 수립 및 시행을 위하여 부동산 거래상황, 외국인 부동산 취득현황, 부동산 가격 동향 등에 관한 정보를 종합적으로 관리하고, 이를 관련 기관ㆍ단체 등에 제공해야 한다.
  \item 광역시장ㆍ도지사는 효율적인 정보의 관리 및 국민편의 증진을 위하여 대통령령으로 정하는 바에 따라 부동산거래의 계약ㆍ신고ㆍ허가ㆍ관리 등의 업무와 관련된 정보체계를 구축ㆍ운영해야 한다.
  \item 국토교통부장관은 정보체계에 구축되어 있는 정보를 수요자에게 제공할 수 있으며, 이 경우 제공하는 정보의 종류와 내용을 제한할 수 없다.
\end{choices}
\end{question}
\answer{2}
